% U7 WS5 -- solubility equilibria, common ion, and a 10-point FRQ. Block 8.
\documentclass{apchem-worksheet}

\apcunit{7}
\apcunittitle{Equilibrium}
\apcdoctitle{Worksheet 5 \textbullet\ Solubility Equilibria \& FRQ}
\apcsubtitle{CED 7.11--7.12 \textbullet\ $K_{sp}$ is a constant; solubility is a position}

\begin{document}
\wsheader[\ce{MX}: $K_{sp} = s^2$ \textbullet\ \ce{MX2} or \ce{M2X}:
$K_{sp} = 4s^3$ \textbullet\ common ion $\Rightarrow$ $Q$ rises above
$K_{sp}$ $\Rightarrow$ less dissolves. \\
$K_{sp}$: \ce{AgCl} $1.6\times10^{-10}$ \textbullet\ \ce{BaSO4}
$1.5\times10^{-9}$ \textbullet\ \ce{CaF2} $4.0\times10^{-11}$
\textbullet\ \ce{PbI2} $1.4\times10^{-8}$ \textbullet\ \ce{Ag2CrO4}
$9.0\times10^{-12}$ \textbullet\ \ce{BaCrO4} $8.5\times10^{-11}$.]

\problem[]{\ced{7.11} Write the dissolution equation and $K_{sp}$
expression:}
\begin{parts}
  \item \ce{AgCl} \answerblank[6.4cm]{$K_{sp} = [\ce{Ag+}][\ce{Cl-}]$}
  \item \ce{CaF2} \answerblank[6.4cm]{$K_{sp} = [\ce{Ca^2+}][\ce{F-}]^2$}
  \item \ce{Ag2CrO4}
        \answerblank[6.4cm]{$K_{sp} = [\ce{Ag+}]^2[\ce{CrO4^2-}]$}
  \item Why does the solid never appear?
        \answerblank[5.4cm]{it is a pure solid}
\end{parts}

\problem[]{\ced{7.11} Calculate the molar solubility in pure water:}
\begin{parts}
  \item \ce{AgCl} \answerblank[3.2cm]{$1.3\times10^{-5}$~M}
  \item \ce{BaSO4} \answerblank[3.2cm]{$3.9\times10^{-5}$~M}
  \item \ce{CaF2} \workspace[2cm]
        \keyonly{\begin{apcanswer}$4s^3 = 4.0\times10^{-11}$;
        $s^3 = 1.0\times10^{-11}$;
        $s = \mathbf{2.2\times10^{-4}}$~M\end{apcanswer}}
  \item \ce{PbI2} \answerblank[3.2cm]{$1.5\times10^{-3}$~M}
\end{parts}

\problem[]{\ced{7.11} \ce{CaF2} has $K_{sp} = 4.0\times10^{-11}$ and
\ce{BaCrO4} has $K_{sp} = 8.5\times10^{-11}$.}
\begin{parts}
  \item Which has the larger $K_{sp}$? \answerblank[2.6cm]{\ce{BaCrO4}}
  \item Calculate both molar solubilities. \workspace[2.6cm]
        \keyonly{\begin{apcanswer}\ce{CaF2}: $s = \sqrt[3]{1.0\times10^{-11}}
        = 2.2\times10^{-4}$~M. \ce{BaCrO4}:
        $s = \sqrt{8.5\times10^{-11}} =
        9.2\times10^{-6}$~M\end{apcanswer}}
  \item Which is actually more soluble, and by what factor?
        \answerblank[5cm]{\ce{CaF2}, about 24 times}
  \item State the rule this illustrates. \answerlines[2]{$K_{sp}$ values may
        be compared directly only for salts that produce the same total
        number of ions, because $s$ enters the expression raised to
        different powers otherwise.}
\end{parts}

\problem[]{\ced{7.12} The common-ion effect.}
\begin{parts}
  \item Calculate the solubility of \ce{AgCl} in \SI{0.10}{\Molar}
        \ce{NaCl}. \workspace[2cm]
        \keyonly{\begin{apcanswer}$[\ce{Cl-}] \approx 0.10$;
        $s = K_{sp}/[\ce{Cl-}] = 1.6\times10^{-10}/0.10 =
        \mathbf{1.6\times10^{-9}}$~M\end{apcanswer}}
  \item By what factor is it suppressed relative to pure water?
        \answerblank[2.6cm]{about 7900}
  \item Explain the effect using Le Ch\^atelier.
        \answerlines[2]{The added \ce{Cl-} is a product of the dissolution
        equilibrium, so adding it shifts the equilibrium left and less
        \ce{AgCl} dissolves.}
  \item Explain the same effect using $Q$.
        \answerlines[3]{Raising $[\ce{Cl-}]$ raises the ion product $Q$
        above $K_{sp}$. A system with $Q > K$ shifts toward the reactant
        side --- here toward undissolved solid --- until $Q$ falls back to
        $K_{sp}$, leaving less dissolved.}
  \item What happened to $K_{sp}$? \answerblank[2.6cm]{nothing}
\end{parts}

\problem[]{\ced{7.12} Calculate the solubility of \ce{CaF2} in
\SI{0.025}{\Molar} \ce{NaF}, and state the assumption you make.
\workspace[2.8cm]
\keyonly{\begin{apcanswer}Assume $[\ce{F-}] \approx 0.025$, since the
dissolution contributes negligibly beside the \ce{NaF}. Then
$4.0\times10^{-11} = s(0.025)^2$, giving
$s = \mathbf{6.4\times10^{-8}}$~M --- about 3400 times less than in pure
water.\end{apcanswer}}}

\problem[]{\ced{7.12} A student says ``the common ion lowers $K_{sp}$,
which is why less dissolves.'' Rewrite the sentence correctly.
\answerlines[3]{The common ion does not change $K_{sp}$, which depends only
on temperature. It shifts the equilibrium \emph{position}: more of the
common ion and correspondingly less of the other ion, with the product of
the two still equal to the same $K_{sp}$.}}

\problem[]{\textbf{FRQ (10 points).} \ced{7.3} \ced{7.7} \ced{7.10}
\ced{7.12} At \SI{700}{\kelvin}, $K = 50.0$ for
$\ce{H2(g) + I2(g) <=> 2HI(g)}$.}
\begin{parts}
  \item Write the equilibrium expression.
        \answerblank[6cm]{$K = [\ce{HI}]^2/([\ce{H2}][\ce{I2}])$}
  \item A flask initially contains \SI{1.00}{\Molar} \ce{H2} and
        \SI{1.00}{\Molar} \ce{I2}. Set up an ICE table and find all three
        equilibrium concentrations. \workspace[3.4cm]
        \keyonly{\begin{apcanswer}Change: $-x$, $-x$, $+2x$.
        $K = (2x)^2/(1.00-x)^2 = 50.0$, so
        $2x/(1.00-x) = 7.07$ and $x = 0.780$.
        $[\ce{H2}] = [\ce{I2}] = \mathbf{0.220}$~M,
        $[\ce{HI}] = \mathbf{1.56}$~M\end{apcanswer}}
  \item Explain why the quadratic formula was not required here.
        \answerlines[2]{Both sides of the expression are perfect squares,
        so taking the square root of each converts it into a linear
        equation in $x$.}
  \item Verify your answer by substituting back into $K$.
        \workspace[1.8cm]
        \keyonly{\begin{apcanswer}$(1.56)^2/(0.220)^2 = 2.43/0.0484
        = \mathbf{50.2}$, agreeing with 50.0 to
        rounding\end{apcanswer}}
  \item A second flask at the same temperature contains
        $[\ce{H2}] = 0.40$, $[\ce{I2}] = 0.30$ and $[\ce{HI}] = 0.40$~M.
        Calculate $Q$ and predict the direction of change.
        \workspace[2.2cm]
        \keyonly{\begin{apcanswer}$Q = (0.40)^2/[(0.40)(0.30)] = 1.3$.
        Since $Q < K$, the reaction proceeds to the \textbf{right} until
        $Q$ rises to 50.0.\end{apcanswer}}
  \item State the effect on $K$ of (i) adding more \ce{H2}, and
        (ii) raising the temperature, given that the forward reaction is
        exothermic. \answerlines[3]{(i) None --- adding \ce{H2} changes
        $Q$, and the system shifts right until $Q$ returns to the same $K$.
        (ii) $K$ decreases, because for an exothermic forward reaction heat
        behaves as a product and adding it pushes the equilibrium toward
        reactants.}
\end{parts}

\begin{apcnote}[Rubric --- score yourself, 10 points]
\textbf{(a) 1 pt} correct expression \textbullet\ \textbf{(b) 3 pts}
1 ICE table with $+2x$ for \ce{HI}, 1 correct algebra to $x = 0.780$,
1 all three concentrations \textbullet\ \textbf{(c) 1 pt} identifies the
perfect square \textbullet\ \textbf{(d) 1 pt} substitution recovering
$\approx 50$ \textbullet\ \textbf{(e) 2 pts} 1 value of $Q$, 1 direction
\emph{with} the $Q < K$ justification --- a bare ``shifts right'' earns 1
\textbullet\ \textbf{(f) 2 pts} 1 states $K$ unchanged for the
concentration change, 1 states $K$ decreases on heating with the exothermic
reason.
\end{apcnote}

\begin{spiralstrip}{Unit 7 \textbullet\ blocks 1--7}
  \item $Q > K$ means the system shifts:
        \answerblank[1.8cm]{left}
  \item Inert gas at constant $V$: \answerblank[2.4cm]{no shift}
  \item $K$ for \ce{MX2} in terms of $s$: \answerblank[1.8cm]{$4s^3$}
  \item Only what alters $K$? \answerblank[2.6cm]{temperature}
  \item Compressing shifts toward: \answerblank[4cm]{fewer moles of gas}
\end{spiralstrip}

\end{document}
